% TODO make hw stylesheet
\documentclass[10pt]{scrartcl}
\usepackage[a4paper, total={6in, 8in}]{geometry}
\usepackage[usenames,dvipsnames,svgnames]{xcolor}
\usepackage[shortlabels]{enumitem}
\usepackage[framemethod=TikZ]{mdframed}
\usepackage{amsmath,amssymb,amsthm}
\usepackage[colorlinks]{hyperref}
\usepackage{multicol}
\usepackage{tabularx}

\hypersetup{
	colorlinks=true,
	linkcolor=blue,
	filecolor=magenta,      
	urlcolor=magenta,
	pdftitle={MAT 322 HW}
}

\usepackage{microtype}
\usepackage{mathtools}
\usepackage[headsepline]{scrlayer-scrpage}
\usepackage{thmtools}
\usepackage[textsize=footnotesize]{todonotes}

\mdfdefinestyle{mdbluebox}{roundcorner=10pt,innerbottommargin=9pt,
	linecolor=blue,backgroundcolor=TealBlue!5,}
\declaretheoremstyle[headfont=\sffamily\bfseries\color{MidnightBlue},
mdframed={style=mdbluebox},]{thmbluebox}

\mdfdefinestyle{mdbloobox}{frametitlefont=\bfseries,innerbottommargin=8pt,
	nobreak=true,backgroundcolor=TealBlue!5,linecolor=blue,}
\declaretheoremstyle[headfont=\bfseries\color{MidnightBlue},
mdframed={style=mdbloobox},headpunct={\\[3pt]},postheadspace=0pt,]{thmbloobox}

\mdfdefinestyle{mdredbox}{frametitlefont=\bfseries,innerbottommargin=8pt,
	nobreak=true,backgroundcolor=Salmon!5,linecolor=RawSienna,}
\declaretheoremstyle[headfont=\bfseries\color{RawSienna},
mdframed={style=mdredbox},headpunct={\\[3pt]},postheadspace=0pt,]{thmredbox}

\mdfdefinestyle{mdgreenbox}{linecolor=ForestGreen,backgroundcolor=ForestGreen!5,
	linewidth=2pt,rightline=false,leftline=true,topline=false,bottomline=false,}
\declaretheoremstyle[headfont=\bfseries\sffamily\color{ForestGreen!70!black},
mdframed={style=mdgreenbox},headpunct={ --- },]{thmgreenbox}

\mdfdefinestyle{mdorangebox}{linecolor=RedOrange,backgroundcolor=RedOrange!5,
	linewidth=2pt,rightline=false,leftline=true,topline=false,bottomline=false,}
\declaretheoremstyle[headfont=\bfseries\sffamily\color{RedOrange!70!black},
mdframed={style=mdorangebox},headpunct={ --- },]{thmorangebox}

\mdfdefinestyle{mdblackbox}{linecolor=black,backgroundcolor=RedViolet!5!gray!5,
	linewidth=3pt,rightline=false,leftline=true,topline=false,bottomline=false,}
\declaretheoremstyle[mdframed={style=mdblackbox}]{thmblackbox}

\declaretheoremstyle[spaceabove=6pt,spacebelow=6pt]{basehead}

\declaretheorem[style=basehead,name=Problem]{problem}
\declaretheorem[style=thmredbox,name=Problem,numbered=no]{problem*}
\declaretheorem[style=thmbloobox,name=Setup]{setup} % for config geo
\declaretheorem[style=thmredbox,name=Example,sibling=problem]{example}
\declaretheorem[style=thmredbox,name=$\bigstar$ Problem,sibling=problem]{niceprob}
\declaretheorem[style=thmbluebox,name=Theorem,sibling=problem]{theorem}
\declaretheorem[style=thmbluebox,name=Corollary,sibling=theorem]{corollary}
\declaretheorem[style=thmbluebox,name=Proposition,sibling=theorem]{prop}
\declaretheorem[style=thmbluebox,name=Lemma,numberwithin=problem]{lemma}
\declaretheorem[style=thmbluebox,name=Theorem,numbered=no]{theorem*}
\declaretheorem[style=thmbluebox,name=Lemma,numbered=no]{lemma*}
\declaretheorem[style=thmgreenbox,name=Claim,numberwithin=problem]{claim}
\declaretheorem[style=thmgreenbox,name=Claim,numbered=no]{claim*}
\declaretheorem[style=thmblackbox,name=Remark,numberwithin=problem]{remark}
\declaretheorem[style=thmblackbox,name=Remark,numbered=no]{remark*}
\declaretheorem[style=thmgreenbox,name=Definition,sibling=theorem]{definition}
\declaretheorem[style=thmgreenbox,name=Definition,numbered=no]{definition*}
\declaretheorem[style=thmorangebox,name=Warning,sibling=theorem]{warning}
\declaretheorem[style=thmorangebox,name=Warning,numbered=no]{warning*}
\declaretheorem[style=thmorangebox,name=Note,numbered=no]{note}
\declaretheorem[style=thmblackbox,name=Exercise,sibling=problem]{exercise}
\declaretheorem[style=thmblackbox,name=Exercise,numbered=no]{exercise*}
\declaretheorem[style=thmblackbox,name=Question,sibling=problem]{question}
\declaretheorem[style=thmblackbox,name=Question,numbered=no]{question*}

\addtolength{\textheight}{3.14cm}
\providecommand{\ol}{\overline}
\providecommand{\ul}{\underline}
\providecommand{\eps}{\varepsilon}
\providecommand{\half}{\frac{1}{2}}
\providecommand{\dang}{\measuredangle} %% Directed angle
\providecommand{\CC}{\mathbb C}
\providecommand{\FF}{\mathbb F}
\providecommand{\NN}{\mathbb N}
\providecommand{\QQ}{\mathbb Q}
\providecommand{\RR}{\mathbb R}
\providecommand{\ZZ}{\mathbb Z}
\providecommand{\dg}{^\circ}
\providecommand{\ii}{\item}
\providecommand{\setdiff}{\smallsetminus}
\providecommand{\alert}[1]{{\sffamily\textbf{\textcolor{blue}{#1}}}}
\providecommand{\inv}{^{-1}}
\providecommand{\mb}[1]{\mathbf{#1}}
\providecommand{\dif}{\;d}
\providecommand{\tensor}{\otimes}

\reversemarginpar

\newenvironment{soln}{\begin{proof}[Solution]}{\end{proof}}
\newenvironment{walkthrough}{\noindent \textbf{Walkthrough.}}{\medskip}
\newlist{walk}{enumerate}{3}
\setlist[walk]{label=\bfseries (\alph*)}

\providecommand{\pow}{\operatorname{Pow}}
\providecommand{\vocab}[1]{{\textbf{\textcolor{ForestGreen}{#1}}}}
\providecommand{\lcm}{\operatorname{lcm}}
\providecommand{\abs}[1]{\left|#1\right|}

\usepackage{asymptote}
\begin{asydef}
	defaultpen(fontsize(10pt));
	unitsize(1cm);
	size(8cm); // set a reasonable default
	usepackage("amsmath");
	usepackage("amssymb");
	settings.tex="pdflatex";
	settings.outformat="pdf";
	// Replacement for olympiad+cse5 which is not standard
	import geometry;
	// recalibrate fill and filldraw for conics
	void filldraw(picture pic = currentpicture, conic g, pen fillpen=defaultpen, pen drawpen=defaultpen)
	{ filldraw(pic, (path) g, fillpen, drawpen); }
	void fill(picture pic = currentpicture, conic g, pen p=defaultpen)
	{ filldraw(pic, (path) g, p); }
	// some geometry
	pair foot(pair P, pair A, pair B) { return foot(triangle(A,B,P).VC); }
	pair orthocenter(pair A, pair B, pair C) { return orthocentercenter(A,B,C); }
	pair centroid(pair A, pair B, pair C) { return (A+B+C)/3; }
	// cse5 abbreviations
	path CP(pair P, pair A) { return circle(P, abs(A-P)); }
	path CR(pair P, real r) { return circle(P, r); }
	pair IP(path p, path q) { return intersectionpoints(p,q)[0]; }
	pair OP(path p, path q) { return intersectionpoints(p,q)[1]; }
	path Line(pair A, pair B, real a=0.6, real b=a) { return (a*(A-B)+A)--(b*(B-A)+B); }
	// cse5 more useful functions
	picture CC() {
		picture p=rotate(0)*currentpicture;
		currentpicture.erase();
		return p;
	}
	pair MP(Label s, pair A, pair B = plain.S, pen p = defaultpen) {
		Label L = s;
		L.s = "$"+s.s+"$";
		label(L, A, B, p);
		return A;
	}
	pair Drawing(Label s = "", pair A, pair B = plain.S, pen p = defaultpen) {
		dot(MP(s, A, B, p), p);
		return A;
	}
	path Drawing(path g, pen p = defaultpen, arrowbar ar = None) {
		draw(g, p, ar);
		return g;
	}
\end{asydef}

\usepackage{mathrsfs}

\ihead{\footnotesize MAT 322 HW09}
\ohead{\footnotesize Last compiled \today}

\title{\vspace{-2em}MAT 322 HW09}
\author{\large Aditya Pahuja}
\date{\large Due 04/29}
\begin{document}
\maketitle

\begin{problem}
	For each of the following functions $T\colon\RR^4\times\RR^4\to\RR$,
	verify that $T$ is a $2$-tensor on $\RR^4$ and express $T$
	in terms of the standard basis of $\mathcal L^2(\RR^4)$.
	\begin{enumerate}[(a)]
		\ii $T(x,y)=x_1y_3+2x_3y_1-3x_2y_4$
		\ii $T(x,y)=7x_4y_2+2x_2y_3+11x_1y_1$
	\end{enumerate}
\end{problem}

\begin{soln}
	We note that the basis elements $e_i^*\otimes e_j^*$ are the tensors
	such that $(e_i^*\otimes e_j^*)(x,y) = x_iy_j$, so for (a) we can write
	\[T = e_1^*\otimes e_3^* + 2e_3^*\otimes e_1^* - 3e_2^*\otimes e_4^*\]
	and for (b) we can write
	\[T = 7e_4^*\otimes e_2^* + 2e_2^*\otimes e_3^* + 11e_1^*\otimes e_1^*.\]
	In either case, we can see that $T$ must be a $2$-tensor on $\RR^4$
	because it is a linear combination of $2$-tensors on $\RR^4$.
\end{soln}

\begin{problem}
	For each of the following functions $\omega\colon\RR^4\times\RR^4\to\RR$,
	verify that $\omega$ is an alternating $2$-tensor on $\RR^4$ and
	express $\omega$ in terms of the standard basis of $\mathcal A^2(\RR^4)$.
	\begin{enumerate}[(a)]
		\ii $\omega(x,y)=x_1y_3-y_1x_3 + 2x_3y_1-2y_3x_1 - 3x_2y_3+3y_2x_3$
		\ii $\omega(x,y)=7x_1y_2-7y_1x_2 + 2x_2y_3-2y_2x_3$
	\end{enumerate}
\end{problem}

\begin{soln}
	We can write $(e_i^*\wedge e_j^*)(x,y) = x_iy_j - x_jy_i$,
	and the standard basis is the set of $e_i^*\wedge e_j^*$ with $i<j$.
	For (a)
	\[\omega = e_1^*\wedge e_3^* + 2e_3^*\wedge e_1^* - 3e_2^*\wedge e_3^*
	= -e_1^*\wedge e_3^* - 3e_2^*\wedge e_3^*,\]
	and for (b)
	\[\omega = 7e_1^*\wedge e_2^* + 2e_2^*\wedge e_3^*.\]
	These are linear combinations of alternating $2$-tensors on $\RR^4$,
	so they must also be alternating $2$-tensors on $\RR^4$,
	since $\mathcal A^2(\RR^4)$ is a linear subspace of $\mathcal L^2(\RR^4)$.
\end{soln}

\begin{problem}
	Consider the differential $1$-forms on $\RR^3$ given by
	\[\omega = xyd x+3dy-yzdz,\quad \eta = xdx-yz^2dy+2xdz.\]
	Verify by direct computation that
	\[d(d\omega) = 0,\quad
	d(\omega\wedge\eta) = d\omega\wedge\eta-\omega\wedge d\eta.\]
\end{problem}

\begin{soln}
	Writing $d\omega = \omega_1dx + \omega_2dy + \omega_3dz$
	(so $\omega_1 = xy$, $\omega_2 = 3$, $\omega_3 = -yz$), we have
	\begin{align*}
		d\omega &=
		d\omega_1\wedge dx + d\omega_2\wedge dy + d\omega_3\wedge dz\\
		&= (ydx + xdy)\wedge dx + 0\wedge dy + (-zdy - ydz)\wedge dz\\
		&= -xdx\wedge dy - zdy\wedge dz.
	\end{align*}
	Then we can compute
	\[d(d\omega) = -dx\wedge dx\wedge dy - dz\wedge dy\wedge dz = 0.\]
	We also see that
	\[\omega\wedge\eta =
	(-xy^2z^2-3x)dx\wedge dy + (2x^2y+xyz)dx\wedge dz+(6x-y^2z^3)dy\wedge dz,\]
	so
	\[d(\omega\wedge\eta) = (-2xy^2z-2x^2-xz+6)dx\wedge dy\wedge dz.\]
	On the other hand,
	\[d\eta = 2dx\wedge dz + 2yzdy\wedge dz,\]
	which gives us
	\[d\omega\wedge\eta - \omega\wedge d\eta =
	(-2x^2-xz)dx\wedge dy\wedge dz - (-6 + 2xy^2z)dx\wedge dy\wedge dz,\]
	which is equal to $(-2xy^2z-2x^2-xz+6)dx\wedge dy\wedge dz
	=d(\omega\wedge\eta)$ as desired.
\end{soln}

\begin{problem}
	Let $\omega$ be a $k$-form on an open set $U\subseteq\RR^n$.
	We say $\omega$ \vocab{vanishes} at $x\in U$ if $\omega_x$
	is the zero tensor.
	\begin{enumerate}[(a)]
		\ii Show that if $\omega$ vanishes for all $x$ in
		a neighborhood of $x_0$, then $d\omega$ vanishes at $x_0$.
		\ii Give an example to show that if $\omega$ vanishes at $x_0$
		then $d\omega$ need not vanish at $x_0$.
	\end{enumerate}
\end{problem}

\begin{soln}
	(a) Let $\omega = \sum_I\omega_Idx_I$ where the sum is over all
	sequences $I = (i_1,\dots,i_k)$ of integers satisfying
	$1\le i_1<\cdots<i_k\le n$, $dx_I$ is shorthand for
	$dx_{i_1}\wedge\cdots\wedge dx_{i_k}$, and the $\omega_I$ are
	smooth functions $U\to\RR$. The vanishing condition means that
	$\omega_I(x) = 0$ for all $x$ in a neighborhood of $x_0$,
	so the derivative of $\omega_I$ is $0$ in a neighborhood of $x_0$
	(in particular the partials are all zero in a neighborhood of $x_0$).
	This means that, for $x$ in this neighborhood,
	\[d\omega_x = \sum_{I}d\omega_I(x)\wedge dx_I = \sum_I 0\wedge dx_I =0.\]

	(b) Take $\omega\colon\RR\to\RR$ to be the $0$-form given by $f(x) = x$
	for all $x$. Then $\omega_0 = 0$, but $d\omega = dx\ne 0$ for all $x$.
\end{soln}

\begin{problem}
	We say that $\omega\in\Omega^k(U)$ is \vocab{exact} if
	there exists $\eta\in\Omega^{k-1}(U)$ such that $\omega = d\eta$,
	and we say that $\omega$ is \vocab{closed} if $d\omega = 0$.
	Note that $d(d\eta) = 0$ implies all exact forms are closed.
	\begin{enumerate}[(a)]
		\ii Show that the $1$-form
		\[\omega = \frac{x}{x^2+y^2}dx + \frac{y}{x^2+y^2}dy\]
		on $\RR^2\setdiff\{(0,0)\}$ is both closed and exact.
		\ii Show that the $1$-form
		\[\omega = \frac{-y}{x^2+y^2}dx + \frac{x}{x^2+y^2}dy\]
		on $\RR^2\setdiff\{(0,0)\}$ is closed.
	\end{enumerate}
\end{problem}

\begin{soln}
	(a) It suffices to show that $\omega$ is exact. Indeed, noting
	$\omega = d\eta$ where $\eta$ is the differential
	$0$-form $\ln(\sqrt{x^2+y^2})$ proves exactness.
	
	(b) This is just a direct derivative computation: we have
	\begin{align*}
		d\omega& = \left(\frac{-y\cdot -2y}{(x^2+y^2)^2}
		- \frac{1}{x^2+y^2}\right)dy\wedge dx + \left(\frac{x\cdot-2x}
		{(x^2+y^2)^2}+\frac{1}{x^2+y^2}\right)dx\wedge dy\\
		&= \left(\frac{y^2-x^2}{(x^2+y^2)^2}\right)(dy\wedge dx+dx\wedge dy)
	\end{align*}
	which obviously equals zero.
\end{soln}























\end{document}